% ============================================================
%  SECCIÓN 2 – Objetivos y Alcance
% ============================================================

\section{Objetivos y Alcance}

\subsection{Objetivo general}

Desarrollar un entorno VR y un esquema de control que permitan entrenar y luego teleoperar un brazo robótico mediante un joystick propio, priorizando seguridad, usabilidad y desempeño.

\subsection{Objetivos específicos}

\begin{enumerate}
    \item Implementar un \textbf{simulador} de brazo robótico electrohidráulico en Unity con interacción XR.
    \item Diseñar el \textbf{mapeo de controles} del joystick (RPi Pico) a los grados de libertad del robot.
    \item Instrumentar \textbf{métricas de entrenamiento}: tiempo por tarea, errores, colisiones, suavidad/jerk.
    \item Diseñar y validar \textbf{modos de teleoperación}: velocidad limitada, asistencia, seguimiento de trayectoria.
    \item Incorporar \textbf{mecanismos de seguridad}: parada de emergencia (\textit{e-stop}), zonas prohibidas, límites de torque/velocidad.
    \item Realizar una \textbf{evaluación} con usuarios: usabilidad (SUS), curva de aprendizaje y transferencia al robot real.
\end{enumerate}

\subsection{Alcance}

\begin{itemize}
    \item Entrenamiento VR de tareas representativas de la aplicación (ej.: \textit{pick \& place} de vasos, posicionamiento preciso dentro del rack).
    \item Teleoperación con joystick RPi Pico (prototipo funcional de hardware y firmware).
    \item Integración de comunicación (serial/UDP/TCP) entre el joystick y Unity.
    \item Capa de seguridad básica en la simulación.
\end{itemize}

\subsection{Fuera de alcance (versión actual)}

\begin{itemize}
    \item Retroalimentación háptica avanzada (vibración proporcional a la fuerza).
    \item \textit{Motion planning} autónomo o visión por computadora.
    \item Integración física directa con el robot electrohidráulico real (la comunicación con el robot real queda como trabajo futuro).
    \item Certificación normativa industrial.
\end{itemize}
